\documentclass[11pt,a4paper]{article}
\usepackage[margin=25mm]{geometry}
\usepackage[T1]{fontenc}
\usepackage{lmodern}
\usepackage{microtype}
\usepackage{booktabs}
\usepackage{amsmath,amssymb,amsthm}
\usepackage{hyperref}
\hypersetup{colorlinks=true,linkcolor=black,urlcolor=blue}
\theoremstyle{definition}
\newtheorem{definition}{Definition}
\newtheorem{assumption}{Assumption}
\newtheorem{property}{Property}
\theoremstyle{remark}
\newtheorem*{remark}{Remark}
\newcommand{\HMAC}{\mathsf{HMAC}}
\newcommand{\Expand}{\mathsf{HKDF\text{-}Expand}}
\newcommand{\SHA}{\mathsf{SHA256}}
\newcommand{\id}{\mathsf{id}}
\newcommand{\cat}{\parallel}
\title{Theoretical Foundation of ZK-SNARK Video Steganography}
\author{VideoLevel technical design note}
\date{19 September 2026 (revised 3 October 2026)}
\begin{document}
\maketitle
\begin{abstract}
This design note formalizes the construction currently implemented by
VideoLevel. A message $m$ and a compressed Groth16 proof $\pi$ are packed into a
payload, wrapped in an HMAC-authenticated frame, whitened with an HMAC
keystream and written into the first trailing-one sign bit of selected CAVLC
residual blocks of H.264 IDR slices. Positions are chosen by a keyed
pseudorandom schedule, so extraction is blind: it needs neither the cover video
nor a sidecar. The note states each security property together with the
assumption it rests on, and lists the properties the construction does not
provide.
\end{abstract}

\paragraph{Update 03/10/2026.}
This revision replaces the earlier description. The key schedule is now
protocol~v2 (HKDF-separated subkeys, per-segment HMAC ordering, frame
whitening); the circuit constrains every hash and secret signal to be boolean;
verification is blind extraction followed by mandatory Groth16 verification.
The former Python pipeline (Python CAVLC parser, position safety filter, chaos
maps for bit/position scrambling, \texttt{positions.json} sidecars, strict mode
that re-parsed the original cover, and the SEC1--SEC10 benchmark suite) has been
removed and is \emph{not} part of the system described here. In particular, no
chaos-map theory applies to the current construction. Streams produced by
protocol~v1 cannot be read by v2.

\section{Notation and trust boundary}
Bytes are big-endian and bit strings are MSB-first. $\mathsf{u}k(x)$ denotes
the $k$-bit big-endian encoding of an integer $x$, $X[a{:}b]$ the byte slice
$a,\ldots,b-1$, and $\HMAC(k,x)$ always means HMAC-SHA256. Let $V$ be a
Baseline-profile H.264 Annex-B elementary stream, $m$ a non-empty byte
string, and $K\in\{0,1\}^{256}$ a 32-byte secret shared by embedder and
receiver. Let $\beta\ge 1$ be the per-segment bit budget
(\texttt{max\_bits\_per\_IDR}) and $L_{\max}$ the configured maximum payload
size (4096 bytes in the native CLI; the frame format allows up to $2^{16}-1$).
Both $\beta$ and $L_{\max}$ are public parameters that embedder and receiver
must agree on.

The embedder outputs a stego stream $V'$ only; no sidecar is produced. The
receiver holds $V'$, $K$, $\beta$, $L_{\max}$ and the Groth16 verification key
$\mathsf{vk}$. The secret $K$ and message $m$ must never be written to job
status records, logs or benchmark artifacts.

\section{Carrier: trailing-one sign bits}
CAVLC codes each $4\times4$ (or DC) residual block as the syntax element
$\mathsf{coeff\_token}$ (the pair $(\mathrm{TotalCoeff},\mathrm{TrailingOnes})$,
VLC table chosen by $nC$), one $\mathsf{trailing\_ones\_sign\_flag}$ per
trailing one, the remaining levels, $\mathsf{total\_zeros}$ and
$\mathsf{run\_before}$.

\begin{definition}[Candidate]
Let $\sigma=(\mathrm{SPS},\mathrm{PPS},\mathrm{IDR})$ be a segment. For every
residual block of category $t\in\{0,1,2,3\}$ (LumaDC, Luma$4{\times}4$,
ChromaDC, ChromaAC) in the IDR slice with $\mathrm{TrailingOnes}\ge1$, the
candidate is the \emph{first} $\mathsf{trailing\_ones\_sign\_flag}$ of that
block in bitstream order. It is identified by
\[
\id(c)=\text{ASCII}\bigl(n:a:t:b:o\bigr),
\]
where $n$ is the NAL index inside the segment, $a$ the macroblock address, $b$
the block index, and $o$ the RBSP bit offset of the sign flag. $C(\sigma)$
denotes the set of candidates of $\sigma$; at most one candidate exists per
block.
\end{definition}

\begin{property}[Length and context preservation]\label{prop:flip}
Flipping a candidate sign changes one coefficient from $\pm1$ to $\mp1$ and
leaves $\mathrm{TotalCoeff}$, $\mathrm{TrailingOnes}$, every codeword length
and every $nC$ value unchanged. Hence the RBSP length, all RBSP bit offsets,
the parse of every later syntax element, and $C(\sigma)$ itself are identical
in cover and stego. (Emulation-prevention bytes are re-inserted when RBSP is
converted back to EBSP, so EBSP byte offsets may shift; identities use RBSP
offsets.)
\end{property}

Property~\ref{prop:flip} follows from the CAVLC syntax and is checked
bit-exactly by the implementation; it is what makes blind extraction possible:
the receiver recomputes $C(\sigma)$ from $V'$. A flip still changes the decoded
picture (a dequantized coefficient change of magnitude $2$ at quantized level,
propagated through intra prediction and deblocking); imperceptibility is an
empirical property measured by decoded PSNR/SSIM, not a theorem. Only IDR
slices carry bits; non-IDR NAL units are copied unchanged.

\section{Key schedule (protocol v2)}
\begin{definition}[Subkeys]
With HKDF-SHA256 (RFC~5869),
\begin{align*}
\mathrm{PRK} &= \HMAC(\texttt{"zkstego-cavlc-v2-salt"},\,K),\\
K_s &= \Expand(\mathrm{PRK},\texttt{"zkstego/cavlc/v2/schedule"},32),\\
K_f &= \Expand(\mathrm{PRK},\texttt{"zkstego/cavlc/v2/frame-tag"},32),\\
K_w &= \Expand(\mathrm{PRK},\texttt{"zkstego/cavlc/v2/whitening"},32).
\end{align*}
\end{definition}

\begin{definition}[Segment schedule]
Let $V'$ be split into segments $\sigma_0,\sigma_1,\ldots$, one per IDR NAL,
each consisting of the latest SPS and PPS followed by that IDR. In segment
$\sigma_j$, order the candidates by the key
$\bigl(\HMAC(K_s,\id(c)),\,\id(c)\bigr)$ ascending (lexicographically on bytes)
and keep the first
\[
q_j=\min\bigl(\beta,\,|C(\sigma_j)|\bigr)
\]
candidates. Concatenating these lists over $j=0,1,\ldots$ gives the global
position sequence $p_0,p_1,\ldots$; frame bit $i$ is written at $p_i$. Bit
indices continue across segments; scheduled positions beyond the frame length
keep their cover sign. Embedding fails if
$\sum_j q_j<|F|$, where $|F|$ is the frame length in bits.
\end{definition}

The identity contains no sign value and $C(\sigma_j)$ is sign-invariant
(Property~\ref{prop:flip}), so a receiver with $K$ and $\beta$ reproduces
$p_0,p_1,\ldots$ exactly from $V'$. A mismatched $\beta$ yields a different
schedule and extraction fails.

\section{Authenticated, whitened frame}\label{sec:frame}
\begin{definition}[Frame and embedded bits]
For a payload $P$ with $|P|\le L_{\max}$,
\[
F=\texttt{0x02}\cat\mathsf{u16}(|P|)\cat P\cat
\HMAC\bigl(K_f,\ \texttt{0x02}\cat\mathsf{u16}(|P|)\cat P\bigr)[0{:}16].
\]
The whitening keystream is
$S=\HMAC(K_w,\mathsf{u64}(0))\cat\HMAC(K_w,\mathsf{u64}(1))\cat\cdots$
and the embedded bit at $p_i$ is
\[
e_i=F_i\oplus S_i,\qquad\text{sign}(p_i)=\begin{cases}+&e_i=0\\-&e_i=1.\end{cases}
\]
\end{definition}
The frame overhead is $1+2+16=19$ bytes, so $|F|=8(|P|+19)$ bits.

\paragraph{Blind extraction.}
(1) Parse the IDR segments of $V'$ and enumerate $C(\sigma_j)$.
(2) Derive $K_s,K_f,K_w$ and the schedule.
(3) Read and de-whiten the first 24 bits; reject unless the version is
\texttt{0x02} and $\mathsf{u16}$ length $\le L_{\max}$.
(4) Read the remaining $8(|P|+16)$ bits, de-whiten, recompute the tag and
compare in constant time; reject on mismatch or if the stream ends early.
With a wrong key the de-whitened version byte is essentially random, so most
wrong keys are rejected at step~(3) and the rest at step~(4).

\section{Payload and zero-knowledge relation}
\begin{definition}[Payload]
$P=\mathsf{u32}(|m|)\cat m\cat\pi$, where $\pi$ is a Groth16 proof over BN254
compressed to 129 bytes (the $x$-coordinates of $A\in G_1$, $B\in G_2$ (two
$\mathbb{F}_p$ elements), $C\in G_1$, 32 bytes each, plus one byte holding
the three $y$-parity bits).
\end{definition}
The length field gives an unambiguous message/proof split. With
$L_{\max}=4096$ the largest message is $4096-4-129=3963$ bytes. Decompression
rejects coordinates $\ge p$ and points not on the curve.

\begin{definition}[Relation]
The circuit \texttt{payload\_verify.circom} (63\,321 constraints, 513 public
signals) defines
\[
\mathcal{R}=\Bigl\{\,(h,c,\ell;\,s)\ :\ c=\SHA(h\cat s),\
h,c,s\in\{0,1\}^{256},\ 0<\ell<10^6\Bigr\},
\]
where the booleanity of every bit of $h$, $c$ and $s$ is enforced by
$x(x-1)=0$ constraints. The prover uses $h=\SHA(m)$, computed \emph{outside}
the circuit, $s=K$ and $\ell=|m|$; the public statement is $(h,c,\ell)$ and
the witness is $s$.
\end{definition}

The booleanity constraints are necessary for soundness: circomlib's
\texttt{Sha256} gadget does not constrain its inputs to bits, so without them
a prover could feed arbitrary field elements into the hash. Note that $\ell$
is only range-checked inside the circuit; the circuit does not relate $\ell$ to
$h$. The link between $\ell$, $h$ and $m$ is established by the verifier
recomputing both from $m$.

\paragraph{Verification.}
The receiver accepts $(m,\pi)$ if and only if all of the following hold:
blind extraction returns a frame with a valid tag; $P$ unpacks into
$\mathsf{u32}(|m|)\cat m\cat\pi$ with a valid 129-byte encoding; and
\[
\mathsf{Groth16.Verify}\bigl(\mathsf{vk},\,(h,c,\ell),\,\pi\bigr)=1
\quad\text{with}\quad h=\SHA(m),\ c=\SHA(h\cat K),\ \ell=|m|
\]
recomputed by the receiver (public signals supplied by the sender are never
used). The HTTP verify job reports \texttt{rejected} with reason
\texttt{payload\_not\_authenticated}, \texttt{malformed\_proof\_payload} or
\texttt{proof\_invalid} when the corresponding step fails.

\section{Security properties}
\begin{assumption}[Standard primitives]\label{ass:prf}
(A1) HMAC-SHA256 is a pseudorandom function (PRF) when keyed with a uniformly
random 32-byte key; consequently HKDF-SHA256 is a secure KDF and
$K_s,K_f,K_w$ are computationally independent pseudorandom keys when $K$ is
uniform. (A2) SHA-256 is collision resistant. (A3) Groth16 is
knowledge-sound and zero-knowledge in the generic-group/algebraic model,
\emph{provided} the setup was generated honestly and its trapdoor destroyed.
\end{assumption}

The following properties hold under Assumption~\ref{ass:prf}; they are
standard reductions and no stronger claim is made.

\begin{property}[Pseudorandom schedule]
For a party without $K$, the ranking of candidates in every segment is
computationally indistinguishable from a uniformly random permutation chosen
independently of the cover (by A1 applied to $K_s$, given distinct
identities). Hence the selected positions are pseudorandom to such a party.
\end{property}

\begin{property}[Frame authenticity]
A party without $K$ that may observe valid frames and makes at most $Q$
verification attempts produces a new frame with a valid tag with probability
at most $\mathrm{Adv}^{\mathrm{prf}}_{\HMAC}+Q\,2^{-128}$, since the tag is a
128-bit truncation of a PRF output. A valid tag therefore shows that the frame
was produced by a holder of $K$ (or of $K_f$) and was not modified.
\end{property}

\begin{property}[No fixed-header bias]
For a single frame per key, $S$ is indistinguishable from uniform by A1 applied
to $K_w$, so the embedded bits $e_i$ are pseudorandom even though $F$ has a
fixed version byte, small length field and structured proof bytes. Selected
signs therefore carry no fixed header pattern.
\end{property}

\begin{property}[Subkey separation]
Disclosure of one subkey (e.g.\ $K_s$) reveals neither $K$ nor the other
subkeys, by A1 applied to $\mathrm{PRK}$ and the HKDF-Expand outputs.
\end{property}

\begin{property}[What the proof shows]
If $\mathsf{Verify}$ accepts, then by A3 the prover knew $s$ with
$\SHA(\SHA(m)\cat s)=c=\SHA(\SHA(m)\cat K)$, and by A2 $s=K$ except with
negligible probability. Groth16 zero-knowledge means $\pi$ reveals nothing
about $s$ beyond this statement.
\end{property}

\subsection*{Properties not provided}
\begin{itemize}
\item \textbf{No binding to video, segment or session.} Neither $F$ nor
$\mathcal{R}$ contains a video hash, segment index, timestamp, nonce or camera
identity. A payload extracted from one stream can be re-embedded into another
stream under the same $K$ and will verify (replay).
\item \textbf{No public verifiability.} The verifier needs $K$ both to extract
(schedule, whitening, tag) and to recompute $c=\SHA(h\cat K)$. Anyone able to
verify can also produce a valid proof, so in the current configuration $\pi$
attests knowledge of the same shared secret that the HMAC tag already attests;
it is not evidence towards third parties.
\item \textbf{No nonce in whitening.} $S$ depends only on $K$. Embedding two
frames under the same $K$ reuses the keystream; the guarantee above is per
single frame per key. Positions remain hidden, which limits but does not
remove this leakage. Whitening is not an IND-CPA encryption of $m$.
\item \textbf{No undetectability proof.} Pseudorandom positions and bits do not
imply that the sign changes are statistically undetectable relative to
natural cover statistics; detectability must be measured.
\item \textbf{Trusted setup.} The current proving/verification keys are a
local educational setup (\texttt{DEMO\_SETUP\_ONLY.txt}); whoever holds the
trapdoor can forge proofs. A3 does not hold for it in a deployment sense.
\item \textbf{Robustness.} The channel is fragile by design: any re-encode,
transcode or bit error in a selected sign breaks the tag.
\end{itemize}

\section{Algorithms}
\subsection*{Embedding}
\begin{enumerate}
\item Validate $|K|=32$, $0<|m|$, $4+|m|+129\le L_{\max}$; load circuit
artifacts.
\item Compute $h=\SHA(m)$, $c=\SHA(h\cat K)$ and $\ell=|m|$. Run Groth16,
$\pi\leftarrow\mathsf{Prove}(\mathsf{pk};(h,c,\ell);K)$, and compress $\pi$ to
129 bytes.
\item Derive $K_s,K_f,K_w$; form $P$ and $F$; for each IDR segment in order,
schedule $q_j$ positions and set $\text{sign}(p_i)$ from $e_i$ until all
$|F|$ bits are placed.
\item Re-insert emulation prevention and emit $V'$; later NAL units are copied.
\end{enumerate}
\subsection*{Verification}
\begin{enumerate}
\item Blind extraction (Section~
ef{sec:frame}) with $K,\beta,L_{\max}$; reject on version,
length, capacity or tag failure.
\item Unpack $P$, decompress $\pi$, reject malformed encodings before pairing
checks.
\item Recompute $(h,c,\ell)$ from $(m,K)$ and run Groth16 verification;
accept only on success.
\end{enumerate}

\section{Capacity and evaluation quantities}
The capacity of a stream is $\sum_j\min(\beta,|C(\sigma_j)|)$ bits and a
message of $|m|$ bytes needs $|F|=8(|m|+4+129+19)=8(|m|+152)$ bits. For
example, an 8-frame CIF all-intra clip at QP~22 has 3660 candidates; a 19-byte
message needs 1368 bits (a 171-byte frame), about 37\% of the candidates.
Evaluations should report candidate counts, scheduled bits, decoded quality
(full-video and minimum modified-frame PSNR, SSIM), verification outcome and
latency per stage separately; measured values are artifact-specific, not
guarantees. The strongest regime is all-intra Baseline CAVLC; with GOP${}>1$,
P-frames referencing modified IDR frames inherit the distortion and need
separate measurement.

\section{Operational surface}
The HTTP service exposes public \texttt{/health} and bearer-token-protected
embed, extract, verify, status and one-time artifact download jobs plus a
WebSocket stream endpoint. The token (\texttt{ZK\_STEGO\_API\_TOKEN}, at least
32 characters) is checked before the body is read; repeated authentication
failures are rate-limited; uploads, concurrency and job queues are bounded.
Status records never contain $K$, $m$ or stack traces. Native processes are run
without a shell and with timeouts.

The signed-manifest and expiring-certificate modules (Ed25519) remain in the
repository as auxiliary trust interfaces but are not part of the embedding or
verification path above. An offline verifier can reject an expired
certificate, but cannot prevent use of a copied key with a modified clock; hard
expiry needs an online lease, trusted time or hardware-enforced key.

\section{Limitations and target design}
The native parser accepts only Baseline CAVLC, progressive 4:2:0, one slice
per frame starting at macroblock~0, no I\_PCM, and a single SPS/PPS id.
Realtime throughput on a physical camera has not met the 30-FPS gate.

\texttt{future\_plan.md} describes a target design that is \textbf{not
implemented}: a relation over a media root, session and segment index,
previous-record digest, policy digest and a camera/registry credential; a
public, sign-independent carrier schedule so that verification no longer needs
a shared secret; segment chaining with a final seal; and removal of HMAC from
the acceptance condition once those components pass their gates. None of the
properties of that design should be attributed to the current construction.
\end{document}
